\chapter*{Sources and further reading}
\addcontentsline{toc}{chapter}{Sources and further reading}
Access checked 17 October 2026. The supplied compiler package was read in full.
SSA, verification, folding, DCE, capability lowering, costs, liveness, source
IDs and exercises were retained and deepened. Its classification of engine
and retrieval operations as freely pure was not retained: our dialect preserves
those observations and checks declared effects against operation semantics.
The independent recursive arithmetic oracle and actual Chapter 39/40
handlers test a more precise contract. Chapter 14 tensor typing,
Chapter 15 transition semantics and Chapter 16 resource conventions were
inspected; this finite IR does not implement every earlier language.

The official MLIR language reference's high-level operation/value model and
side-effect rationale's categorization/modeling sections were inspected.
The TVM primary arXiv abstract was inspected for compiler organization only.
This reference implements neither framework, target code generation nor
hardware experiments. Chapter 36 integration and production serialization remain open.
\begingroup
\renewcommand{\bibname}{References}
\let\chapter\section
\begin{thebibliography}{9}
\bibitem{evod41:mlir} MLIR contributors.
\emph{MLIR Language Reference}. Official documentation, high-level structure
and operations/values inspected 17 October 2026.
\url{https://mlir.llvm.org/docs/LangRef/}.
\bibitem{evod41:effects} MLIR contributors.
\emph{Side Effects \& Speculation}. Official rationale, overview,
categorization and modeling sections inspected 17 October 2026.
\url{https://mlir.llvm.org/docs/Rationale/SideEffectsAndSpeculation/}.
\bibitem{evod41:tvm} Tianqi Chen et al.
\emph{TVM: An Automated End-to-End Optimizing Compiler for Deep Learning}.
2018. Primary abstract inspected; no benchmark reproduction.
\url{https://arxiv.org/abs/1802.04799}.
\end{thebibliography}
\endgroup
